\section{What Is Predictable in Sub-Second Link State?}
\label{sec:predictable}
Before we build a forecaster, we ask which properties of the RTT carry over from the input window to the forecast horizon. The statistics in this section use the measured pings of the training sessions, with each session centered on its own mean so that differences between sessions do not appear as memory. Confidence intervals come from a bootstrap over sessions.

\begin{figure*}[t]
  \centering
  % fr_volatility.tex - data from modeling/experiments/fr_volatility.py (results_final/analysis/volatility/).
% Autocorrelation on TRAINING sessions of the RTT level, the RTT change dRTT and the change size |dRTT|,
% lags 0.2-10 s, pairs inside one session only, each session centred on its own mean, real pings only.
% Shaded bands: 95% session-bootstrap intervals. Insets zoom into the boxed region (lags up to 3 s, near zero)
% with the same range in every inset. Starlink DL has no inset: its 1 s sawtooth is visible at full scale.
% Open black circles at 0.2 s: the same ACF after shuffling the real pings inside every session (memoryless null).
% Where a curve starts on its circle, the first lag is mechanical (one spike = two large changes), not memory.
% Needs: \input{figs/data/plotstyle.tex}, pgfplots groupplots, tikz calc. No fillbetween library.
\begin{tikzpicture}
\tikzset{
  volZero/.style={black!65, line width=0.4pt},
  volBox/.style={black!70, line width=0.5pt},
  volArrow/.style={->, >=stealth, black!70, line width=0.5pt, shorten >=1pt},
}
\pgfplotsset{
  volL/.style={cPers, densely dashed, thick, mark=none},
  volD/.style={cResid, thick, mark=triangle*, mark size=1.5pt, mark indices={1,5,10,15,20,25,30,35,40,45,50}},
  volA/.style={cFR, very thick, mark=square*, mark size=1.1pt, mark indices={1,5,10,15,20,25,30,35,40,45,50}},
  volBandD/.style={draw=none, fill=cResid, fill opacity=0.28, forget plot},
  volBandA/.style={draw=none, fill=cFR, fill opacity=0.28, forget plot},
  volInset/.style={scale only axis, width=2.45cm, height=1.5cm, xmin=0.1, xmax=3.1, ymin=-0.15, ymax=0.3,
    xtick={0.2,1,2,3}, xticklabels={0.2,1,2,3}, ytick={-0.1,0,0.1,0.2,0.3}, yticklabels={,0,,0.2,},
    tick label style={font=\tiny}, axis background/.style={fill=white}, grid=major, grid style={gray!20},
    xticklabel style={fill=white, inner sep=0.5pt}, yticklabel style={fill=white, inner sep=0.5pt},
    clip=true},
  volInD/.style={cResid, thick, mark=triangle*, mark size=1.0pt},
  volInA/.style={cFR, thick, mark=square*, mark size=0.75pt},
  volNullMark/.style={only marks, mark=o, mark size=2.2pt, black, line width=0.6pt},
  volNull/.style={volNullMark, restrict x to domain=0:0.3, forget plot},
}
\begin{groupplot}[group style={group size=3 by 2, horizontal sep=0.32cm, vertical sep=0.9cm,
    y descriptions at=edge left, x descriptions at=edge bottom},
  frbase, scale only axis, width=5.05cm, height=3.3cm, xmin=0, xmax=10, ymin=-0.58, ymax=1.0,
  xtick={0,2,4,6,8,10}, ytick={-0.5,0,0.5,1}, xlabel={lag (s)}, ylabel={autocorrelation}]
% ---------------- (a) OneWeb DL
\nextgroupplot[title={(a) OneWeb DL}, legend to name=legvol, legend columns=-1,
  legend style={draw=black!60, fill=white, font=\scriptsize, /tikz/every even column/.append style={column sep=0.3cm}}]
\addplot[volBandD] table[x=x, y=d] {figs/data/fr_volatility_band_oneweb_dl.dat} -- cycle;
\addplot[volBandA] table[x=x, y=absd] {figs/data/fr_volatility_band_oneweb_dl.dat} -- cycle;
\draw[volZero] (axis cs:0,0) -- (axis cs:10,0);
\addplot[volL] table[x=lag_s, y=rtt] {figs/data/fr_volatility_acf_oneweb_dl.dat}; \addlegendentry{RTT level}
\addplot[volD] table[x=lag_s, y=d] {figs/data/fr_volatility_acf_oneweb_dl.dat}; \addlegendentry{RTT change $\Delta$RTT}
\addplot[volA] table[x=lag_s, y=absd] {figs/data/fr_volatility_acf_oneweb_dl.dat}; \addlegendentry{change size $|\Delta\mathrm{RTT}|$}
\addplot[volNull] table[x=lag_s, y=dnull] {figs/data/fr_volatility_acf_oneweb_dl.dat};
\addplot[volNull] table[x=lag_s, y=absdnull] {figs/data/fr_volatility_acf_oneweb_dl.dat};
\addlegendimage{area legend, draw=none, fill=black!22}\addlegendentry{95\% session bootstrap}
\addlegendimage{volNullMark}\addlegendentry{memoryless null}
\draw[volBox] (axis cs:0.1,-0.15) rectangle (axis cs:3.1,0.3);
\coordinate (zA) at (axis cs:3.1,0.3);
% ---------------- (b) Starlink DL (no inset)
\nextgroupplot[title={(b) Starlink DL}]
\addplot[volBandD] table[x=x, y=d] {figs/data/fr_volatility_band_starlink_dl.dat} -- cycle;
\addplot[volBandA] table[x=x, y=absd] {figs/data/fr_volatility_band_starlink_dl.dat} -- cycle;
\draw[volZero] (axis cs:0,0) -- (axis cs:10,0);
\addplot[volL] table[x=lag_s, y=rtt] {figs/data/fr_volatility_acf_starlink_dl.dat};
\addplot[volD] table[x=lag_s, y=d] {figs/data/fr_volatility_acf_starlink_dl.dat};
\addplot[volA] table[x=lag_s, y=absd] {figs/data/fr_volatility_acf_starlink_dl.dat};
\addplot[volNull] table[x=lag_s, y=dnull] {figs/data/fr_volatility_acf_starlink_dl.dat};
\addplot[volNull] table[x=lag_s, y=absdnull] {figs/data/fr_volatility_acf_starlink_dl.dat};
\node[font=\scriptsize, fill=white, inner sep=1.5pt, anchor=north east] at (axis cs:9.8,0.97) {peaks every 1 s};
% ---------------- (c) 5G DL
\nextgroupplot[title={(c) 5G DL}]
\addplot[volBandD] table[x=x, y=d] {figs/data/fr_volatility_band_quectel_dl.dat} -- cycle;
\addplot[volBandA] table[x=x, y=absd] {figs/data/fr_volatility_band_quectel_dl.dat} -- cycle;
\draw[volZero] (axis cs:0,0) -- (axis cs:10,0);
\addplot[volL] table[x=lag_s, y=rtt] {figs/data/fr_volatility_acf_quectel_dl.dat};
\addplot[volD] table[x=lag_s, y=d] {figs/data/fr_volatility_acf_quectel_dl.dat};
\addplot[volA] table[x=lag_s, y=absd] {figs/data/fr_volatility_acf_quectel_dl.dat};
\addplot[volNull] table[x=lag_s, y=dnull] {figs/data/fr_volatility_acf_quectel_dl.dat};
\addplot[volNull] table[x=lag_s, y=absdnull] {figs/data/fr_volatility_acf_quectel_dl.dat};
\draw[volBox] (axis cs:0.1,-0.15) rectangle (axis cs:3.1,0.3);
\coordinate (zC) at (axis cs:3.1,0.3);
% ---------------- (d) OneWeb UL
\nextgroupplot[title={(d) OneWeb UL}]
\addplot[volBandD] table[x=x, y=d] {figs/data/fr_volatility_band_oneweb_ul.dat} -- cycle;
\addplot[volBandA] table[x=x, y=absd] {figs/data/fr_volatility_band_oneweb_ul.dat} -- cycle;
\draw[volZero] (axis cs:0,0) -- (axis cs:10,0);
\addplot[volL] table[x=lag_s, y=rtt] {figs/data/fr_volatility_acf_oneweb_ul.dat};
\addplot[volD] table[x=lag_s, y=d] {figs/data/fr_volatility_acf_oneweb_ul.dat};
\addplot[volA] table[x=lag_s, y=absd] {figs/data/fr_volatility_acf_oneweb_ul.dat};
\addplot[volNull] table[x=lag_s, y=dnull] {figs/data/fr_volatility_acf_oneweb_ul.dat};
\addplot[volNull] table[x=lag_s, y=absdnull] {figs/data/fr_volatility_acf_oneweb_ul.dat};
\draw[volBox] (axis cs:0.1,-0.15) rectangle (axis cs:3.1,0.3);
\coordinate (zD) at (axis cs:3.1,0.3);
% ---------------- (e) Starlink UL
\nextgroupplot[title={(e) Starlink UL}]
\addplot[volBandD] table[x=x, y=d] {figs/data/fr_volatility_band_starlink_ul.dat} -- cycle;
\addplot[volBandA] table[x=x, y=absd] {figs/data/fr_volatility_band_starlink_ul.dat} -- cycle;
\draw[volZero] (axis cs:0,0) -- (axis cs:10,0);
\addplot[volL] table[x=lag_s, y=rtt] {figs/data/fr_volatility_acf_starlink_ul.dat};
\addplot[volD] table[x=lag_s, y=d] {figs/data/fr_volatility_acf_starlink_ul.dat};
\addplot[volA] table[x=lag_s, y=absd] {figs/data/fr_volatility_acf_starlink_ul.dat};
\addplot[volNull] table[x=lag_s, y=dnull] {figs/data/fr_volatility_acf_starlink_ul.dat};
\addplot[volNull] table[x=lag_s, y=absdnull] {figs/data/fr_volatility_acf_starlink_ul.dat};
\draw[volBox] (axis cs:0.1,-0.15) rectangle (axis cs:3.1,0.3);
\coordinate (zE) at (axis cs:3.1,0.3);
% ---------------- (f) 5G UL
\nextgroupplot[title={(f) 5G UL}]
\addplot[volBandD] table[x=x, y=d] {figs/data/fr_volatility_band_quectel_ul.dat} -- cycle;
\addplot[volBandA] table[x=x, y=absd] {figs/data/fr_volatility_band_quectel_ul.dat} -- cycle;
\draw[volZero] (axis cs:0,0) -- (axis cs:10,0);
\addplot[volL] table[x=lag_s, y=rtt] {figs/data/fr_volatility_acf_quectel_ul.dat};
\addplot[volD] table[x=lag_s, y=d] {figs/data/fr_volatility_acf_quectel_ul.dat};
\addplot[volA] table[x=lag_s, y=absd] {figs/data/fr_volatility_acf_quectel_ul.dat};
\addplot[volNull] table[x=lag_s, y=dnull] {figs/data/fr_volatility_acf_quectel_ul.dat};
\addplot[volNull] table[x=lag_s, y=absdnull] {figs/data/fr_volatility_acf_quectel_ul.dat};
\draw[volBox] (axis cs:0.1,-0.15) rectangle (axis cs:3.1,0.3);
\coordinate (zF) at (axis cs:3.1,0.3);
\end{groupplot}
% ---------------- insets: lags up to 3 s near zero, same range in every inset
\foreach \g/\c/\ds in {A/c1r1/oneweb_dl, C/c3r1/quectel_dl, D/c1r2/oneweb_ul, E/c2r2/starlink_ul, F/c3r2/quectel_ul}{
  \begin{axis}[volInset, at={($(group \c.north east)+(-0.1cm,-0.12cm)$)}, anchor=north east, name=ins\g]
  \addplot[volBandD] table[x=x, y=d] {figs/data/fr_volatility_band_\ds.dat} -- cycle;
  \addplot[volBandA] table[x=x, y=absd] {figs/data/fr_volatility_band_\ds.dat} -- cycle;
  \draw[volZero] (axis cs:0,0) -- (axis cs:3.2,0);
  \addplot[volInD] table[x=lag_s, y=d] {figs/data/fr_volatility_acf_\ds.dat};
  \addplot[volInA] table[x=lag_s, y=absd] {figs/data/fr_volatility_acf_\ds.dat};
  \end{axis}
  \draw[volArrow] (z\g) -- (ins\g.south west);
}
\node at ($(group c2r1.north)+(0,0.8cm)$) {\pgfplotslegendfromname{legvol}};
\end{tikzpicture}

  \caption{Autocorrelation of the RTT level, of its 200\,ms change $\Delta$RTT, and of the size of the change $|\Delta\mathrm{RTT}|$ in the training sessions (shaded: 95\% session-bootstrap interval). Circles: value of a memoryless series at 0.2\,s. A curve that starts on its circle has no memory at the first lag. Insets: lags of 0.2 to 3\,s. Starlink DL peaks every second because of a load-driven sawtooth.}
  \label{fig:vol}
\end{figure*}

\subsection{The Next Change Is Hard, the Level Reverts}
Fig.~\ref{fig:vol} shows the autocorrelation of the RTT level, of its change $\Delta$RTT between consecutive rows, and of the size of the change $|\Delta\mathrm{RTT}|$. At the first lag, the change of a memoryless series has an autocorrelation of $-0.5$ (circles), because two consecutive changes share one sample. From 0.4\,s to 10\,s the autocorrelation of $\Delta$RTT stays within $\pm$\VolDmaxAbs{} on five links. Past changes therefore say almost nothing about the next change. This does not make the direction of the next change unpredictable, because the RTT returns toward its recent level. Expecting a return to the mean of the 6\,s input window gets the sign of the 1\,s change right in \VolSignRevLo\% to \VolSignRevHi\% of the test windows (balanced accuracy, 50\% for a constant guess), and the learned forecasters of Section~\ref{sec:results} reach \VolSignNNLo\% to \VolSignNNHi\%. Point forecasts thus beat persistence mainly by not repeating the last, noisy ping. Scaled by the error of persistence, a common reference for forecast accuracy~\cite{hyndman2006another}, their MAE is \PersRatioLo{} to \PersRatioHi{} times lower.

\subsection{The Spread Depends on the Link State}
The size of the change has a short memory. The autocorrelation of $|\Delta\mathrm{RTT}|$ at 1\,s is \VolAbsdLo{} to \VolAbsdHi{} on OneWeb and 5G DL, and on the five links without a sawtooth it stays positive for 1.4 to 5.8\,s. Inside a session, this memory is too weak to forecast the next seconds. On OneWeb and 5G, the standard deviation (SD) of the 6\,s input does not predict the error of a 1\,s forecast within a session (rank correlation \VolWithinLo{} to \VolWithinHi), but it does between sessions (\VolBetweenLo{} to \VolBetweenHi). Each session runs in a noise regime set by its load and radio conditions, and a few seconds of input reveal the regime. Starlink DL also has dynamics within a session. Under the saturating downlink test, its RTT climbs from about 35\,ms to about 200\,ms in steps of about 50\,ms and drops back once per second, which produces the peaks of Fig.~\ref{fig:vol}(b). The RTT autocorrelation at 1\,s is \VolSLDLlevel{} and exceeds \VolSLDLlevelMin{} in every one of the \VolSLDLnSess{} training sessions. The sawtooth switches on and off within a session, as in Fig.~\ref{fig:trace}, and the volatility of the last 6\,s predicts that of the next 3\,s (within-session rank correlation \VolSLDLrv). We attribute the sawtooth to queue build-up and drain under the saturating load, since the uplink test does not show it.

\subsection{Errors Are Skewed Upward}
The RTT also moves up further than it moves down. Queues and retransmissions add delay, while nothing removes delay below the propagation floor. For the 1\,s forecasts of the previous model, the 97.5\% quantile of the error on the calibration sessions is \SkewQUL{} times the size of the 2.5\% quantile on 5G UL and \SkewOWDL{} times on OneWeb DL. A symmetric interval around a point forecast wastes width below the forecast.

\subsection{Consequence for Uncertainty}
An interval of constant width around a point forecast ignores both properties. Calibrated to cover 95\% of all windows, it covers \CovConstCalmLo{} to \CovConstCalmHi{} of the calmest tenth of windows, ranked by the input SD, and \CovConstVolLo{} to \CovConstVolHi{} of the most volatile tenth. The calm windows are covered more often on every link. On Starlink DL the interval is \WidConstSLDL\,ms wide at 1\,s ahead in every window, while a quiet window rarely moves by more than a few milliseconds. By Proposition~\ref{prop:optimal}, the optimal interval follows the conditional quantiles. A forecaster should therefore learn the spread and its skew from the input. Because the regime differs between sessions, its intervals must be calibrated on sessions that it has not seen.
